% Auto-generated complete chapter from 09_ara_ai_tools_and_secondary_projects.md
\chapter{09 — ARA, AI Tools and Secondary Projects Strategy}
\label{complete:root_011}

\begin{sourcemeta}
\textbf{Fuente:} \href{file:///E:/Laboral/09_ara_ai_tools_and_secondary_projects.md}{09\_ara\_ai\_tools\_and\_secondary\_projects.md}\\
\textbf{Seccion:} root\\
\textbf{Estado:} canonical\\
\textbf{Rol:} Modulo principal\\
\textbf{Lineas originales:} 840\\
\textbf{Ultima modificacion:} 2026-06-11 13:34\\
\textbf{Deduplicacion conservadora:} 0 bloques repetidos omitidos; 0 caracteres repetidos no reimpresos.
\end{sourcemeta}

\section*{Resumen de orientacion}
This file defines how Alexander Woodcock Salomón should use ARA Framework, ai-news-aggregator, FitApp-Free, the hyperrealistic Blender portrait and minor repositories inside the broader international job-search strategy.

\section*{Contenido completo depurado}
\addcontentsline{toc}{section}{Contenido completo depurado}




\subsection{0. Purpose of this section}




This file defines how Alexander Woodcock Salomón should use ARA Framework, ai-news-aggregator, FitApp-Free, the hyperrealistic Blender portrait and minor repositories inside the broader international job-search strategy.




The goal is not to expand every project. The goal is to decide which projects support the primary positioning, which projects should remain secondary, which require cleanup before being public, and which should be omitted to avoid profile dilution.




The primary positioning remains:




\begin{quote}
Real-Time 3D Developer / Unity Technical Artist focused on interactive technical visualization, CAD-to-realtime optimization, and WebGL deployment.
\end{quote}




Secondary projects must support that positioning. They must not replace it.




\medskip\hrule\medskip




\subsection{1. Source-of-truth assumptions}




\begin{center}
\scriptsize
\begin{longtable}{>{\raggedright\arraybackslash}p{0.455\linewidth} >{\raggedright\arraybackslash}p{0.455\linewidth}}
\toprule
Item & Current status \\
\midrule
\endfirsthead
\toprule
Item & Current status \\
\midrule
\endhead
Thesis & Complete, pending defense \\
Flagship project & TwinSight X500 \\
Primary market route & Unity / Real-Time 3D / Technical Visualization \\
Secondary route & Tools / Pipeline / Python automation \\
AI route & Supporting route, not primary yet \\
Art route & Supporting evidence, not main positioning \\
Full-stack route & Fallback only \\
Fitness/product route & Low priority for current job search \\
\bottomrule
\end{longtable}
\end{center}




The practical implication is simple: secondary projects are useful only if they strengthen one of these claims:




\begin{itemize}
\item technical systems thinking;
\item automation mindset;
\item ability to build tools;
\item ability to document technical workflows;
\item 3D production literacy;
\item ability to finish and present a coherent product.
\end{itemize}




If a project does not support those claims, it should not be highlighted.




\medskip\hrule\medskip




\subsection{2. Project hierarchy}




\begin{center}
\scriptsize
\begin{longtable}{>{\raggedright\arraybackslash}p{0.305\linewidth} >{\raggedright\arraybackslash}p{0.305\linewidth} >{\raggedright\arraybackslash}p{0.305\linewidth}}
\toprule
Tier & Project & Public priority \\
\midrule
\endfirsthead
\toprule
Tier & Project & Public priority \\
\midrule
\endhead
1 & TwinSight X500 & Mandatory \\
2 & ARA Framework & Conditional \\
2 & Hyperrealistic Blender portrait & Conditional \\
3 & ai-news-aggregator & Low / conditional \\
3 & FitApp-Free & Low \\
4 & deus-ex-machina & Audit before showing \\
4 & skills-introduction-to-github & Do not highlight \\
4 & old web/client projects & Only if relevant and clean \\
\bottomrule
\end{longtable}
\end{center}




The portfolio should not look like a random archive of experiments. It should look like a technical professional profile with a main axis and controlled secondary proof.




\medskip\hrule\medskip




\subsection{3. ARA Framework}




\subsubsection{3.1 Current positioning}




ARA Framework is an academic research automation system / multi-agent research workflow prototype. It uses Python, LangGraph, LangChain, Groq LLaMA 3.3-70B, Redis, optional Supabase, Semantic Scholar API, Playwright MCP, MarkItDown MCP, PDF processing, web scraping, checkpointing, budget management and Markdown report generation.




It should be described as:




\begin{quote}
AI-assisted research automation prototype for structured academic and technical report generation.
\end{quote}




It should not be described as:




\begin{itemize}
\item enterprise AI platform;
\item production-ready research operating system;
\item autonomous research agent comparable to commercial products;
\item ML/AI engineering proof;
\item peer-reviewed research infrastructure;
\item finished SaaS.
\end{itemize}




\subsubsection{3.2 Strategic value}




ARA is useful because it shows:




\begin{itemize}
\item Python systems thinking;
\item workflow decomposition;
\item agent orchestration awareness;
\item API integration;
\item retrieval/research workflow design;
\item automation mindset;
\item technical documentation;
\item ability to use AI beyond casual prompting.
\end{itemize}




ARA is not yet strong enough to be the main employment hook unless it has:




\begin{itemize}
\item working demo;
\item clear README;
\item installation instructions;
\item example output;
\item architecture diagram;
\item limitations;
\item clean environment setup;
\item no secrets in repo;
\item controlled claims.
\end{itemize}




\subsubsection{3.3 Relevance by target role}




\begin{center}
\scriptsize
\begin{longtable}{>{\raggedright\arraybackslash}p{0.455\linewidth} >{\raggedright\arraybackslash}p{0.455\linewidth}}
\toprule
Role family & ARA relevance \\
\midrule
\endfirsthead
\toprule
Role family & ARA relevance \\
\midrule
\endhead
Real-Time 3D / Interactive 3D & Low / indirect \\
Unity Technical Artist & Medium if framed as tooling mindset \\
Technical Visualization & Medium if framed as documentation automation \\
Unity WebGL & Low \\
Tools / Pipeline Developer & High \\
Python Automation Developer & High \\
LLM Application Developer & Medium / conditional \\
Full-stack Developer & Medium if UI/API exists \\
Game Technical Artist & Low / indirect \\
\bottomrule
\end{longtable}
\end{center}




\subsubsection{3.4 Public/private decision}




ARA should become public only after a cleanup pass.




Recommended status now:




\begin{quote}
Keep private or semi-private until it has a recruiter-safe README and one reproducible demo.
\end{quote}




Minimum public release requirements:




\begin{itemize}
\item no API keys;
\item no private paths;
\item no broken scripts in root;
\item \texttt{README.md} explaining purpose, architecture, setup and limitations;
\item \texttt{.env.example};
\item example input and output;
\item architecture diagram;
\item screenshots or terminal demo;
\item short demo video or GIF;
\item license decision;
\item explicit “prototype” label.
\end{itemize}




\subsubsection{3.5 README structure for ARA}




Recommended README sections:




\begin{mdcode}
Bloque de codigo: text
# ARA Framework

AI-assisted research automation prototype for structured academic and technical report generation.

## Problem
Manual academic research workflows are repetitive, fragmented and hard to reproduce.

## What it does
- niche analysis
- literature discovery
- technical architecture draft
- implementation research
- report synthesis

## Architecture
- Niche Analyst
- Literature Researcher
- Technical Architect
- Implementation Specialist
- Content Synthesizer

## Tech stack
Python, LangGraph, LangChain, Redis, Semantic Scholar API, Playwright MCP, MarkItDown MCP, Groq LLaMA.

## Current status
Prototype. Functional core under cleanup. Not a production system.

## Demo
Input example, output example, screenshots.

## Limitations
- depends on external APIs
- outputs require human review
- source quality varies
- not a substitute for academic judgment

## Future work
- stronger evaluation
- better source validation
- UI layer
- cost controls
\end{mdcode}




\subsubsection{3.6 Claims allowed}




Allowed:




\begin{itemize}
\item “Built a Python prototype for structured research automation.”
\item “Designed a multi-agent workflow using LangGraph.”
\item “Integrated external research and document-processing tools.”
\item “Implemented checkpointing, report generation and modular agent responsibilities.”
\item “Used AI assistance while retaining architecture, debugging and integration responsibility.”
\end{itemize}




Not allowed unless evidence exists:




\begin{itemize}
\item “Production-ready.”
\item “Autonomous research platform.”
\item “Enterprise-grade.”
\item “AI engineer.”
\item “Machine learning engineer.”
\item “Deployed SaaS.”
\item “Scientifically validated.”
\end{itemize}




\subsubsection{3.7 ARA action plan}




\begin{center}
\scriptsize
\begin{longtable}{>{\raggedright\arraybackslash}p{0.305\linewidth} >{\raggedright\arraybackslash}p{0.305\linewidth} >{\raggedright\arraybackslash}p{0.305\linewidth}}
\toprule
Priority & Task & Purpose \\
\midrule
\endfirsthead
\toprule
Priority & Task & Purpose \\
\midrule
\endhead
1 & Remove secrets and private config & Security \\
2 & Verify reproducible run & Credibility \\
3 & Add example input/output & Recruiter comprehension \\
4 & Add architecture diagram & Technical proof \\
5 & Add limitations & Trust \\
6 & Add demo video/GIF & Fast review \\
7 & Decide public/private & Exposure control \\
\bottomrule
\end{longtable}
\end{center}




ARA should not receive major development time before TwinSight portfolio assets are published.




\medskip\hrule\medskip




\subsection{4. ai-news-aggregator}




\subsubsection{4.1 Current positioning}




ai-news-aggregator is a private full-stack AI news aggregation and analysis platform. It reportedly uses FastAPI, React, PostgreSQL, Redis, Python backend and an analysis pipeline.




It is useful only if it can be demonstrated as a working system.




Current recommended status:




\begin{quote}
Keep private until a minimal demo works.
\end{quote}




\subsubsection{4.2 Strategic value}




Potential value:




\begin{itemize}
\item FastAPI proof;
\item React proof;
\item PostgreSQL / Redis proof;
\item product logic;
\item data ingestion;
\item AI summarization pipeline;
\item deployment potential.
\end{itemize}




Current risk:




\begin{itemize}
\item it dilutes the 3D/technical-art positioning;
\item it makes the profile look like unfocused full-stack/AI exploration;
\item if incomplete, it weakens credibility;
\item if public but broken, it is worse than invisible.
\end{itemize}




\subsubsection{4.3 Role relevance}




\begin{center}
\scriptsize
\begin{longtable}{>{\raggedright\arraybackslash}p{0.455\linewidth} >{\raggedright\arraybackslash}p{0.455\linewidth}}
\toprule
Role family & Relevance \\
\midrule
\endfirsthead
\toprule
Role family & Relevance \\
\midrule
\endhead
Real-Time 3D / Interactive 3D & Low \\
Unity Technical Artist & Low \\
Technical Visualization & Low / indirect \\
Tools / Pipeline Developer & Medium \\
Python Automation Developer & Medium / high \\
LLM Application Developer & Medium \\
Full-stack Developer & High \\
\bottomrule
\end{longtable}
\end{center}




\subsubsection{4.4 Publication decision}




Do not pin this repo unless it has:




\begin{itemize}
\item working local setup;
\item live demo or screenshots;
\item seeded sample data;
\item clean README;
\item architecture diagram;
\item no broken dependencies;
\item no private API keys;
\item clear explanation of what AI does and does not do.
\end{itemize}




\subsubsection{4.5 Minimum viable public version}




The minimum version does not need to be impressive. It needs to be reliable.




Required:




\begin{itemize}
\item article ingestion from one or two sources;
\item dashboard with list and detail view;
\item summary or categorization pipeline;
\item backend API documented;
\item local setup command;
\item screenshots;
\item known limitations.
\end{itemize}




Optional:




\begin{itemize}
\item user accounts;
\item saved topics;
\item PostgreSQL deployment;
\item Docker compose;
\item evaluation metrics;
\item LangGraph integration.
\end{itemize}




\subsubsection{4.6 Recommendation}




Do not prioritize ai-news-aggregator in the first 30 days of job-search preparation.




Use it later only if applying to:




\begin{itemize}
\item Python automation roles;
\item LLM application roles;
\item full-stack roles;
\item internal tooling roles.
\end{itemize}




It should not appear above TwinSight or ARA.




\medskip\hrule\medskip




\subsection{5. FitApp-Free}




\subsubsection{5.1 Current positioning}




FitApp-Free is a fitness / calisthenics / multi-discipline app in development. It focuses on progressions, injury considerations, routine tracking, workout tracking, weight progression, repetitions, user progress and exercise programming logic.




\subsubsection{5.2 Strategic value}




FitApp may show:




\begin{itemize}
\item product thinking;
\item domain modeling;
\item routine/progression logic;
\item UX/system design;
\item persistence/data structures;
\item long-term product initiative.
\end{itemize}




But it does not strongly support the main profile.




\subsubsection{5.3 Risk}




FitApp is the highest dilution risk among the named projects because it shifts the reader toward:




\begin{itemize}
\item generic app development;
\item fitness product management;
\item personal-product hobby;
\item incomplete startup idea.
\end{itemize}




For the current strategy, this is mostly noise.




\subsubsection{5.4 Recommended status}




\begin{quote}
Keep private or omit from recruiter-facing materials for now.
\end{quote}




Use only if:




\begin{itemize}
\item applying to full-stack/product roles;
\item it has a clean demo;
\item it shows strong system design;
\item it does not displace TwinSight.
\end{itemize}




\subsubsection{5.5 Minimum public threshold}




Before showing FitApp publicly, it needs:




\begin{itemize}
\item working user flow;
\item exercise/progression model;
\item routine builder;
\item tracking UI;
\item screenshots;
\item README explaining system design;
\item no incomplete placeholder state.
\end{itemize}




\subsubsection{5.6 Recommendation}




Do not spend significant job-search time on FitApp during the 0–6 month employment sprint.




\medskip\hrule\medskip




\subsection{6. Hyperrealistic Blender portrait}




\subsubsection{6.1 Current positioning}




The Blender portrait was made 100\% in Blender following CG Cookie HUMAN. It demonstrates character pipeline, topology, materials, grooming, lighting, rendering and production discipline.




It is useful only if presented as a technical breakdown, not as a generic render.




\subsubsection{6.2 Strategic value}




The portrait can support:




\begin{itemize}
\item Blender fluency;
\item topology awareness;
\item material setup;
\item grooming;
\item lighting and rendering;
\item ability to follow a high-end character pipeline;
\item visual quality standards.
\end{itemize}




It is less relevant for:




\begin{itemize}
\item Unity WebGL;
\item technical visualization;
\item CAD optimization;
\item simulation;
\item tools development.
\end{itemize}




\subsubsection{6.3 Recommended framing}




Recommended title:




\begin{quote}
Hyperrealistic Character Pipeline Study — Blender
\end{quote}




Not recommended:




\begin{quote}
Senior Character Artist Portfolio Piece
\end{quote}




\subsubsection{6.4 ArtStation structure}




Recommended ArtStation page:




\begin{enumerate}
\item Final render.
\item Clay render.
\item Wireframe / topology.
\item UV layout.
\item Material breakdown.
\item Grooming breakdown.
\item Lighting setup.
\item Reference / course context.
\item What was learned.
\item Limitations.
\end{enumerate}




\subsubsection{6.5 Disclosure}




Because the work followed a course, the page should say:




\begin{quote}
Character pipeline study completed in Blender following CG Cookie HUMAN training material. Final execution, scene setup, grooming/material work and presentation prepared independently as a portfolio learning piece.
\end{quote}




This avoids implying it was a fully original character art commission or production asset.




\subsubsection{6.6 Role relevance}




\begin{center}
\scriptsize
\begin{longtable}{>{\raggedright\arraybackslash}p{0.455\linewidth} >{\raggedright\arraybackslash}p{0.455\linewidth}}
\toprule
Role family & Relevance \\
\midrule
\endfirsthead
\toprule
Role family & Relevance \\
\midrule
\endhead
Unity Technical Artist & Medium \\
Real-Time 3D Developer & Low / medium \\
Technical Visualization & Low \\
3D Artist / Generalist & High \\
Game Technical Artist & Medium \\
Tools / Pipeline & Low \\
\bottomrule
\end{longtable}
\end{center}




\subsubsection{6.7 Recommendation}




Include it in ArtStation and portfolio as a secondary piece after TwinSight.




It should not be the first portfolio item.




\medskip\hrule\medskip




\subsection{7. deus-ex-machina}




\subsubsection{7.1 Current status}




Public project. Exact current quality must be audited before use.




\subsubsection{7.2 Potential value}




Depending on content, it may show:




\begin{itemize}
\item older game development work;
\item GameMaker/GML experience;
\item design iteration;
\item early programming interest;
\item narrative/gameplay prototyping.
\end{itemize}




\subsubsection{7.3 Risk}




If immature or messy, it may signal:




\begin{itemize}
\item outdated tools;
\item low polish;
\item unfinished habits;
\item lack of focus.
\end{itemize}




\subsubsection{7.4 Recommendation}




Do not pin by default.




Use only if:




\begin{itemize}
\item it has screenshots or gameplay video;
\item it has a clean README;
\item it supports a game-dev application;
\item it is framed as “early independent game prototype,” not as a flagship technical project.
\end{itemize}




\medskip\hrule\medskip




\subsection{8. skills-introduction-to-github}




\subsubsection{8.1 Current status}




Public training / tutorial repository.




\subsubsection{8.2 Recommendation}




Do not highlight. Do not pin.




It can remain public, but it should not be part of the recruiter-facing portfolio.




Reason:




\begin{itemize}
\item it does not demonstrate differentiated capability;
\item it may make the GitHub profile look beginner-level if pinned;
\item it consumes attention without adding evidence.
\end{itemize}




\medskip\hrule\medskip




\subsection{9. Other lower-priority repositories}




Named repositories include:




\begin{itemize}
\item Chocolateria;
\item Bosque de Sombras Manuales;
\item alexander-fullstack-portfolio;
\item alexander-3d-portfolio;
\item delarge95.github.io;
\item 3DWeb;
\item job-search-strategy;
\item ludex\_framework;
\item WoodSleiman;
\item WebGL\_tesis.
\end{itemize}




These require a repository-level audit before publication decisions.




\subsubsection{9.1 Audit criteria}




Each repository should be scored against:




\begin{center}
\scriptsize
\begin{longtable}{>{\raggedright\arraybackslash}p{0.455\linewidth} >{\raggedright\arraybackslash}p{0.455\linewidth}}
\toprule
Criterion & Pass condition \\
\midrule
\endfirsthead
\toprule
Criterion & Pass condition \\
\midrule
\endhead
Relevance & Supports target role \\
Completeness & Has functional output \\
README & Explains purpose and setup \\
Visual proof & Screenshots/video/demo \\
Code hygiene & No broken or private content \\
Recruiter value & Understandable in 60 seconds \\
Dilution risk & Does not confuse positioning \\
\bottomrule
\end{longtable}
\end{center}




\subsubsection{9.2 Default decisions}




\begin{center}
\scriptsize
\begin{longtable}{>{\raggedright\arraybackslash}p{0.455\linewidth} >{\raggedright\arraybackslash}p{0.455\linewidth}}
\toprule
Repository type & Default action \\
\midrule
\endfirsthead
\toprule
Repository type & Default action \\
\midrule
\endhead
Old portfolio repos & Merge/replace with current portfolio \\
Client web projects & Show only if professional and permission-safe \\
Job-search repo & Keep private \\
Manual/document repos & Omit unless relevant \\
Duplicate thesis repos & Consolidate or archive \\
Experimental prototypes & Keep private unless polished \\
\bottomrule
\end{longtable}
\end{center}




\subsubsection{9.3 GitHub visibility rule}




Public does not mean featured.




Pinned repositories should be fewer and stronger. A recruiter should see a coherent technical profile within the first 30–60 seconds.




\medskip\hrule\medskip




\subsection{10. Recommended GitHub pinned repositories}




\subsubsection{10.1 Immediate target state}




\begin{center}
\scriptsize
\begin{longtable}{>{\raggedright\arraybackslash}p{0.305\linewidth} >{\raggedright\arraybackslash}p{0.305\linewidth} >{\raggedright\arraybackslash}p{0.305\linewidth}}
\toprule
Pin order & Repository & Status \\
\midrule
\endfirsthead
\toprule
Pin order & Repository & Status \\
\midrule
\endhead
1 & WebGL-Thesis-Proposal / TwinSight X500 & Mandatory \\
2 & ARA Framework & Conditional after cleanup \\
3 & Portfolio site & Mandatory once active \\
4 & Blender portrait breakdown repo or page & Optional \\
5 & ai-news-aggregator & Conditional \\
6 & one small polished tool/demo & Optional \\
\bottomrule
\end{longtable}
\end{center}




\subsubsection{10.2 Repos to avoid pinning}




\begin{itemize}
\item skills-introduction-to-github;
\item incomplete FitApp;
\item old duplicated thesis repos;
\item unmaintained tests;
\item unclear client projects;
\item any repo with missing README or broken setup.
\end{itemize}




\medskip\hrule\medskip




\subsection{11. Secondary project case-study formats}




\subsubsection{11.1 ARA Framework case study}




Recommended structure:




\begin{enumerate}
\item Problem.
\item System goal.
\item Workflow architecture.
\item Agent roles.
\item Tech stack.
\item Input/output example.
\item Integration challenges.
\item Human validation layer.
\item Limitations.
\item Future work.
\end{enumerate}




Key message:




\begin{quote}
ARA demonstrates automation and systems design, not production AI authority.
\end{quote}




\subsubsection{11.2 ai-news-aggregator case study}




Recommended structure:




\begin{enumerate}
\item Problem.
\item Data ingestion.
\item Backend architecture.
\item Frontend UI.
\item Analysis pipeline.
\item Storage model.
\item Deployment status.
\item Limitations.
\item Future work.
\end{enumerate}




Key message:




\begin{quote}
Shows full-stack and AI-product experimentation; secondary to real-time 3D.
\end{quote}




\subsubsection{11.3 FitApp-Free case study}




Recommended structure if ever public:




\begin{enumerate}
\item Problem.
\item Exercise progression model.
\item Routine logic.
\item Tracking architecture.
\item UI flow.
\item Data model.
\item Limitations.
\item Future work.
\end{enumerate}




Key message:




\begin{quote}
Shows product thinking, not target-role proof.
\end{quote}




\subsubsection{11.4 Blender portrait case study}




Recommended structure:




\begin{enumerate}
\item Final result.
\item Goal.
\item Reference/training context.
\item Modeling/topology.
\item Materials.
\item Grooming.
\item Lighting/rendering.
\item Lessons learned.
\item Limitations.
\end{enumerate}




Key message:




\begin{quote}
Demonstrates 3D production literacy and visual standards.
\end{quote}




\medskip\hrule\medskip




\subsection{12. AI-assisted development disclosure across projects}




\subsubsection{12.1 General rule}




Do not hide AI assistance if directly asked. Do not over-emphasize it unnecessarily.




The safest formula:




\begin{quote}
I used AI assistance as part of the development workflow for ideation, code scaffolding, debugging and documentation support. I was responsible for architecture, integration, testing, technical decisions and final implementation.
\end{quote}




\subsubsection{12.2 Project-specific disclosure}




\begin{center}
\scriptsize
\begin{longtable}{>{\raggedright\arraybackslash}p{0.455\linewidth} >{\raggedright\arraybackslash}p{0.455\linewidth}}
\toprule
Project & Disclosure level \\
\midrule
\endfirsthead
\toprule
Project & Disclosure level \\
\midrule
\endhead
TwinSight & Moderate if asked \\
ARA & Explicit; AI is part of the system \\
ai-news-aggregator & Moderate \\
FitApp & Low unless asked \\
Blender portrait & Mention course context, not AI unless used \\
\bottomrule
\end{longtable}
\end{center}




\subsubsection{12.3 Claims to avoid}




Avoid:




\begin{itemize}
\item “I coded everything manually” if AI was involved;
\item “AI built it” because that undermines ownership;
\item “fully autonomous” unless evaluated;
\item “production AI system” unless deployed and monitored;
\item “research-grade” unless validated.
\end{itemize}




\medskip\hrule\medskip




\subsection{13. Project-to-role mapping}




\begin{center}
\scriptsize
\begin{longtable}{>{\raggedright\arraybackslash}p{0.305\linewidth} >{\raggedright\arraybackslash}p{0.305\linewidth} >{\raggedright\arraybackslash}p{0.305\linewidth}}
\toprule
Project & Best-supported roles & Weak roles \\
\midrule
\endfirsthead
\toprule
Project & Best-supported roles & Weak roles \\
\midrule
\endhead
TwinSight X500 & Real-Time 3D, Unity WebGL, Technical Visualization & AI Tools \\
ARA Framework & Python Automation, Tools, LLM Apps & Unity WebGL \\
ai-news-aggregator & Full-stack, AI Apps & Technical Art \\
FitApp-Free & Product / Full-stack & Technical Visualization \\
Blender portrait & 3D Generalist, Technical Art support & AI / Full-stack \\
deus-ex-machina & Game prototype & Technical Visualization \\
\bottomrule
\end{longtable}
\end{center}




\medskip\hrule\medskip




\subsection{14. Short-term project time allocation}




For the 0–6 month job-search sprint:




\begin{center}
\scriptsize
\begin{longtable}{>{\raggedright\arraybackslash}p{0.455\linewidth} >{\raggedright\arraybackslash}p{0.455\linewidth}}
\toprule
Area & Recommended effort \\
\midrule
\endfirsthead
\toprule
Area & Recommended effort \\
\midrule
\endhead
TwinSight case study & 45\% \\
Portfolio site & 20\% \\
LinkedIn/CV/GitHub cleanup & 15\% \\
ARA cleanup & 10\% \\
Blender portrait presentation & 5\% \\
Other projects & 5\% \\
\bottomrule
\end{longtable}
\end{center}




This allocation is intentionally strict. The main risk is over-dispersing into too many unfinished projects.




\medskip\hrule\medskip




\subsection{15. Decision gates}




\subsubsection{Gate 1 — Before public release}




A project can be public-facing only if it has:




\begin{itemize}
\item clear purpose;
\item clean README;
\item screenshots or video;
\item no secrets;
\item no broken critical path;
\item accurate status label;
\item no exaggerated claims.
\end{itemize}




\subsubsection{Gate 2 — Before pinning}




A project can be pinned only if it supports the target role within 60 seconds.




\subsubsection{Gate 3 — Before including in CV}




A project can enter the CV only if it supports the specific job target.




Examples:




\begin{itemize}
\item Unity Technical Artist CV: TwinSight + Blender portrait + ARA lightly.
\item Python Automation CV: ARA + ai-news-aggregator + TwinSight lightly.
\item Full-stack CV: ai-news-aggregator + FitApp + ARA.
\item Technical Visualization CV: TwinSight first; ARA only if relevant.
\end{itemize}




\medskip\hrule\medskip




\subsection{16. Recommended final project portfolio order}




Recommended public order:




\begin{enumerate}
\item TwinSight X500 — flagship case study.
\item ARA Framework — tools/automation case study after cleanup.
\item Hyperrealistic Blender Portrait — technical 3D breakdown.
\item ai-news-aggregator — optional technical appendix.
\item FitApp-Free — omit unless targeting product/full-stack.
\item Older repos — archive or omit unless cleaned.
\end{enumerate}




\medskip\hrule\medskip




\subsection{17. Risks and mitigations}




\begin{center}
\scriptsize
\begin{longtable}{>{\raggedright\arraybackslash}p{0.455\linewidth} >{\raggedright\arraybackslash}p{0.455\linewidth}}
\toprule
Risk & Mitigation \\
\midrule
\endfirsthead
\toprule
Risk & Mitigation \\
\midrule
\endhead
Profile dilution & Keep TwinSight first \\
Overclaiming AI skill & Label ARA as prototype \\
Broken public repos & Audit before pinning \\
Too many unfinished projects & Limit active cleanup to 2 projects \\
Course-project ambiguity & Disclose course context \\
Recruiter confusion & Use role-specific portfolio variants \\
\bottomrule
\end{longtable}
\end{center}




\medskip\hrule\medskip




\subsection{18. Immediate action checklist}




\subsubsection{18.1 ARA}




\begin{itemize}
\item [ ] Confirm current runnable state.
\item [ ] Remove secrets.
\item [ ] Add \texttt{.env.example}.
\item [ ] Add architecture diagram.
\item [ ] Add example input/output.
\item [ ] Add limitations.
\item [ ] Add 60-second demo.
\item [ ] Decide public/private.
\end{itemize}




\subsubsection{18.2 ai-news-aggregator}




\begin{itemize}
\item [ ] Verify whether it still runs.
\item [ ] Fix minimum broken dependencies.
\item [ ] Capture screenshots.
\item [ ] Decide if worth publishing.
\item [ ] Keep private if not demonstrable.
\end{itemize}




\subsubsection{18.3 FitApp-Free}




\begin{itemize}
\item [ ] Do not prioritize now.
\item [ ] Keep private.
\item [ ] Revisit only for full-stack/product route.
\end{itemize}




\subsubsection{18.4 Blender portrait}




\begin{itemize}
\item [ ] Export final render.
\item [ ] Export clay render.
\item [ ] Export topology screenshots.
\item [ ] Export UV/material/groom breakdown.
\item [ ] Publish as ArtStation technical study.
\end{itemize}




\subsubsection{18.5 GitHub}




\begin{itemize}
\item [ ] Pin TwinSight first.
\item [ ] Remove beginner/tutorial repos from pinned area.
\item [ ] Add README to every visible repo.
\item [ ] Archive duplicate thesis/portfolio repos.
\item [ ] Keep job-search strategy private.
\end{itemize}




\medskip\hrule\medskip




\subsection{19. Final decision}




ARA should be used as a secondary differentiator, not as the main identity. It strengthens a tools/pipeline/automation narrative but should not compete with TwinSight as the flagship.




ai-news-aggregator and FitApp-Free should not receive major attention during the initial job-search sprint unless the strategy pivots toward full-stack or AI application roles.




The Blender portrait is worth presenting, but only as a technical 3D breakdown and clearly below TwinSight.




The portfolio must communicate one core message:




\begin{quote}
Alexander builds interactive real-time 3D technical systems, and also has supporting automation/tooling and 3D production skills.
\end{quote}



